\par\medskip\noindent
\textbf{\textit{\SGARefLink{sga3:xix:corollary:2:6}{Corollary 2.6}.}}---
\textit{Let \(G\) be an \(S\)-group scheme that is affine and smooth
over \(S\), with connected fibers.  The set of points \(s\in S\) for
which \(G_{\bar s}\) is reductive \textup{(resp. semisimple)} is an
open subset \(U\) of \(S\), and the functions}
\[
  s\longmapsto\operatorname{rgred}(G_{\bar s}/\bar s),
  \qquad
  s\longmapsto\operatorname{rgss}(G_{\bar s}/\bar s)
\]
\textit{are locally constant on \(U\).}
\SGARefTarget{sga3:xix:corollary:2:6}

\par\medskip\noindent
\textbf{\textit{\SGARefLink{sga3:xix:definition:2:7}{Definition 2.7}.}}---
\textit{An \(S\)-group \textup{(that is, an \(S\)-group scheme)}
\(G\) is called \emph{reductive} \textup{(resp. \emph{semisimple})}
if it is affine and smooth over \(S\), with connected reductive
\textup{(resp. semisimple)} geometric fibers.}
\SGARefTarget{sga3:xix:definition:2:7}

For an \(S\)-group scheme \(G\), the property of being reductive
(resp. semisimple) is stable under base change and local for the fpqc
topology.

\par\medskip\noindent
\textbf{2.8.}---
Let \(G\) be a reductive \(S\)-group.  For every torus
(resp. maximal torus) \(Q\) of \(G\),
\[
  Z(Q)=\underline{\operatorname{Centr}}_G(Q)
\]
is reductive (resp. is equal to \(Q\)).  This follows at once from \SGARefLink{sga3:xix:paragraph:2-2}{2.2}
and 1.6.  Applying \SGARefLink{sga3:xix:theorem:2-5}{2.5} to
\(\underline{\operatorname{Centr}}_G(Q)\), we deduce that \(Q\) is
contained, locally for the étale topology, in a maximal torus.
\SGARefTarget{sga3:xix:paragraph:2:8}

\par\medskip\noindent
\textbf{\textit{Remark 2.9.}}---
\textit{Using, as usual, the technique of EGA IV\textsubscript{3},
§8, the finite-presentation hypotheses, and Theorem 2.5, one sees
that if \(G\) is a reductive group over \(S\), there is an open cover
\(\{U_i\}\) of \(S\) such that each \(G_{U_i}\) is obtained by base
change from a reductive group over a noetherian ring \textup{(in fact,
over a finitely generated \(\mathbb Z\)-algebra)}.  Likewise, if \(G\)
has a split maximal torus \(T\), one may further suppose that
\(T_{U_i}\) comes from a split maximal torus of the preceding
group, \ldots}
\SGARefTarget{sga3:xix:remark:2:9}

\SGASection{3}{Roots and root systems of reductive group schemes}
\SGARefTarget{sga3:xix:section:3}

